\documentclass[../report.tex]{subfiles}

\begin{document}

\subsection{Linux - Complete Pentest (Mesa)}
\subsubsection{Conditions}
The penetration test was carried out in the same laboratory environment as the first assignment.
Attacks on the target system had to be carried out with the same Kali Linux machine as in the previous assignment as well.

The target machine was set up in the Institute for Cybersecurity’s laboratory again, whose network was made accessible via VPN. It was a Linux machine running Debian, hosting multiple web-applications. It was reachable via SSH, allowing anyone with the right credentials inside the network to access the machine. 

The participants were given a brief about the backstory of the fictitious company \enquote{Mesa}, for which the machine was set up in the experiment. The assignment was to penetration test the company infrastructure. The scope was defined to only include the server with the IP \textit{192.168.61.65}. Other machines, as well as examination of private customer data, were explicitly out of scope.

\subsubsection{Goals}
The overall goal of this assignment was to collect all publicly available information that could aid in an attack on the company's infrastructure and use it to find all relevant vulnerabilities in the system.

There were multiple ways to obtain information that could lead to the compromise of credentials that needed to be discovered, which could be used to log into the organizations web application and the machine running the web servers. Ways of logging in without credential should be tested as well.
Avenues to create reverse shells were present and needed to be discovered, as well as escalating limited privileges of already compromised accounts. 

A vulnerability in an outdated piece of software was present as well, for which a public exploit had to be used.
\end{document}